\documentclass[11pt]{article}
\usepackage[margin=1in]{geometry}
\usepackage{amsmath,amssymb,amsthm,mathtools}
\usepackage[T1]{fontenc}
\usepackage{lmodern}
\usepackage{microtype}
\usepackage{hyperref}
\usepackage{enumitem}
\usepackage{booktabs}

\hypersetup{colorlinks=true,linkcolor=blue,citecolor=blue,urlcolor=blue}

\newtheorem{prop}{Proposition}
\newtheorem{remark}{Remark}

\newcommand{\Hh}{\mathbb{H}}
\newcommand{\Ss}{\mathbb{S}}
\newcommand{\Rr}{\mathbb{R}}
\newcommand{\Vol}{\operatorname{Vol}}
\newcommand{\conv}{\operatorname{conv}}
\newcommand{\Om}{\Omega}
\newcommand{\diag}{\operatorname{diag}}

\title{Notes on hyperbolic $5$- and $6$-simplex volumes\\
formulas, reductions, and a practical code path}
\author{}
\date{}

\begin{document}
\maketitle

\begin{abstract}
This note records a workable extension of the earlier $\Hh^4$ discussion to compact simplices
in $\Hh^5$ and $\Hh^6$.  In dimension $5$ the natural tool is Schlaefli's differential formula,
which reduces the problem to hyperbolic tetrahedral $3$-faces.  In dimension $6$ the even-dimensional
Poincare-Schlaefli-Hopf reduction gives an absolute volume formula in terms of the $21$ dihedral
angles, the $35$ spherical tetrahedral links of the $2$-faces, and the $7$ spherical $5$-simplex
links of the vertices.  We also explain the numerical strategy used in the companion Python module:
QMC evaluation of Klein-model simplex integrals, intrinsic extraction of lower-dimensional Klein
coordinates from ambient hyperboloid Gram data, and gnomonic evaluation of spherical link volumes.
\end{abstract}

\section{General simplex data from Klein vertices}

Let
\[
\Delta=[x_0,\dots,x_n]\subset B^n=\{x\in \Rr^n:\|x\|<1\}
\]
be a compact hyperbolic $n$-simplex in the Klein ball.  Its hyperboloid lifts are
\[
 h_i=\frac{(1,x_i)}{\sqrt{1-\|x_i\|^2}}\in \Hh^n\subset \Rr^{n,1},
 \qquad
 \langle u,v\rangle=-u_0v_0+u_1v_1+\cdots+u_nv_n.
\]
The unnormalized Klein matrix is
\[
Q=(Q_{ij}),\qquad Q_{ij}=1-x_i\cdot x_j.
\]
Since the true vertex Gram matrix is
\[
G_{ij}=-\langle h_i,h_j\rangle=
\frac{1-x_i\cdot x_j}{\sqrt{(1-\|x_i\|^2)(1-\|x_j\|^2)}},
\]
normalization cancels in the inverse, and if $R=Q^{-1}$ then the angle Gram matrix is
\begin{equation}
\label{eq:angle-gram-general}
A_{ii}=1,
\qquad
A_{ij}=-\frac{R_{ij}}{\sqrt{(-R_{ii})(-R_{jj})}}
\quad (i\neq j).
\end{equation}
Hence the dihedral angle between the facets opposite $x_i$ and $x_j$ is
\begin{equation}
\label{eq:dihedral-general}
\theta_{ij}=\arccos(-A_{ij}).
\end{equation}
Formula \eqref{eq:angle-gram-general} is the same linear-algebra package already used in the
$\Hh^4$ note; it is just dimension-independent.

For numerical cross-checks it is useful to remember the direct Klein integral
\begin{equation}
\label{eq:klein-direct}
\Vol_n(\Delta)=\int_{\Delta_E}\frac{dx}{(1-\|x\|^2)^{(n+1)/2}},
\end{equation}
where $\Delta_E=\conv(x_0,\dots,x_n)\subset \Rr^n$ is the ordinary Euclidean simplex in the Klein
ball.  The companion code evaluates \eqref{eq:klein-direct} by Sobol QMC sampling on the standard
simplex.

\section{The $\Hh^5$ case}

\subsection*{Schlaefli differential formula}

For a smoothly varying family of compact hyperbolic $5$-polytopes,
Schlaefli's differential formula reads
\begin{equation}
\label{eq:schlafli-5}
 d\Vol_5(\Delta)
 =-
 \frac{1}{4}
 \sum_{F^3}
 \Vol_3(F^3)\,d\theta_F,
\end{equation}
where the sum runs over the codimension-$2$ faces, i.e. the $3$-dimensional tetrahedral faces.
For a $5$-simplex there are $\binom{6}{2}=15$ such faces.  This is the correct replacement for the
$\Hh^4$ ``absolute'' reduction: in odd dimension one no longer gets volume directly from the
alternating angle sum, so one integrates Schlaefli instead \cite{Kohno,Heckman}.

If the simplex is
\[
\Delta=[x_0,\dots,x_5]\subset B^5,
\]
set
\[
\Delta_t=[tx_0,\dots,tx_5],\qquad 0<t\le 1.
\]
Then $\Delta_t$ stays compact and collapses to the origin as $t\to 0$.  Choosing some small
$t_0\in(0,1)$, one has
\begin{equation}
\label{eq:h5-radial}
\Vol_5(\Delta)
=
\Vol_5(\Delta_{t_0})
-
\frac14
\int_{t_0}^{1}
\sum_{i<j}
\Vol_3\bigl(F_{ij}(t)\bigr)
\,d\theta_{ij}(t),
\end{equation}
where $F_{ij}(t)$ is the tetrahedral face spanned by the four vertices with indices different from
$i,j$.

\subsection*{Extracting the intrinsic $\Hh^3$ face geometry}

The code does not implement the hyperbolic tetrahedron dilogarithm formula explicitly,
although that would also be possible using Murakami--Yano/Ushijima.
Instead it extracts intrinsic Klein coordinates for each tetrahedral face and applies the direct
$3$-dimensional Klein integral.

Suppose $u_0,\dots,u_3\in B^5$ are the four Klein vertices of a tetrahedral face.  Lift them to the
hyperboloid as above, and form the vertex Gram matrix
\[
G^{\mathrm{face}}_{ab}=-\langle h_a,h_b\rangle,\qquad 0\le a,b\le 3.
\]
Equivalently, the Minkowski inner-product matrix of the lifted vertices is
\[
M_{ab}=\langle h_a,h_b\rangle.
\]
Because the face is intrinsically $3$-dimensional, $M$ has signature $(1,3)$.  Numerically one
factors it as
\[
M=C^{T}J_{3,1}C,
\qquad
J_{3,1}=\diag(-1,1,1,1),
\]
so that the columns of $C$ are hyperboloid coordinates
\[
 y_a=(y_{a,0},y_{a,1},y_{a,2},y_{a,3})\in \Hh^3.
\]
Then the intrinsic Klein coordinates of the face are simply
\begin{equation}
\label{eq:intrinsic-klein-face}
 z_a=\frac{(y_{a,1},y_{a,2},y_{a,3})}{y_{a,0}}\in B^3.
\end{equation}
Finally,
\begin{equation}
\label{eq:intrinsic-face-volume}
\Vol_3(F)=\int_{\conv(z_0,z_1,z_2,z_3)}\frac{dz}{(1-\|z\|^2)^2}.
\end{equation}
Combining \eqref{eq:h5-radial} and \eqref{eq:intrinsic-face-volume} gives a completely explicit
numerical pipeline from the original five-dimensional Klein vertices.

\begin{remark}
At the theoretical level, hyperbolic tetrahedron volumes are already dilogarithmic, and Kellerhals
showed that in hyperbolic $5$-space the general odd-dimensional picture reaches trilogarithms,
with doubly asymptotic $5$-orthoschemes providing a fundamental building block
\cite{Kellerhals5,MY,Ushijima}.  The code in this note stays closer to the raw simplex data and uses
QMC for the lower-dimensional face integrals instead.
\end{remark}

\section{The $\Hh^6$ case}

\subsection*{The even-dimensional Hopf reduction}

Let $w_k$ denote the sum of the normalized $(6-k)$-fold interior angles at the $k$-faces of a
$6$-simplex.  Thus $w_6=1$; $w_4$ is built from the $21$ codimension-$2$ dihedral angles; $w_2$
comes from the $35$ spherical tetrahedral links of the $2$-faces; and $w_0$ comes from the $7$
spherical $5$-simplex vertex links.

The Peschl relations imply
\begin{equation}
\label{eq:W-six}
W(\Delta^6)=w_0-\frac12 w_2+w_4-\frac{17}{4},
\end{equation}
while the hyperbolic Poincare--Schlaefli--Hopf theorem gives
\begin{equation}
\label{eq:Hopf-six}
W(\Delta^6)=-\frac{2}{\Om_6}\Vol_6(\Delta^6),
\qquad
\Om_6=\Vol(\Ss^6)=\frac{16\pi^3}{15}.
\end{equation}
Therefore
\begin{equation}
\label{eq:h6-main}
\boxed{
\Vol_6(\Delta)
=
\frac{8\pi^3}{15}
\left(
\frac{17}{4}-w_4+\frac12 w_2-w_0
\right).
}
\end{equation}
This is the $6$-dimensional analogue of the $\Hh^4$ formula, except that now the unresolved term is
an $\Ss^5$-volume rather than an $\Ss^3$-volume \cite{KellerhalsZehrt,Heckman}.

\subsection*{Writing the three pieces from the angle Gram matrix}

Let
\[
\Delta=[x_0,\dots,x_6]\subset B^6,
\]
and let $A$ be its angle Gram matrix from \eqref{eq:angle-gram-general}.  Then:

\paragraph{The $w_4$ term.}
The codimension-$2$ faces correspond to pairs $\{i,j\}$, so
\begin{equation}
\label{eq:w4}
 w_4=\frac{1}{2\pi}\sum_{0\le i<j\le 6}\theta_{ij}.
\end{equation}

\paragraph{The $w_2$ term.}
A $2$-face is the intersection of four facets, hence corresponds to a $4$-subset
$I\subset\{0,\dots,6\}$.  The spherical tetrahedral link at that face has angle Gram matrix equal to
the principal submatrix $A_I$.  Therefore
\begin{equation}
\label{eq:w2}
 w_2=\frac{1}{2\pi^2}\sum_{|I|=4}\Vol_{\Ss^3}(L_I),
\qquad
\text{with }\operatorname{Gram}(L_I)=A_I.
\end{equation}
Each $\Vol_{\Ss^3}(L_I)$ is computable by Murakami's explicit logarithm/dilogarithm formula.
This is exactly the same special-function kernel that already appeared in the $\Hh^4$ vertex-link
term \cite{MurakamiSpherical}.

\paragraph{The $w_0$ term.}
The vertex link at $x_i$ is the spherical $5$-simplex obtained by deleting row and column $i$:
\begin{equation}
\label{eq:w0}
 w_0=\frac{1}{\pi^3}\sum_{i=0}^{6}\Vol_{\Ss^5}(L_i),
 \qquad
 \text{with }\operatorname{Gram}(L_i)=A^{(i)}.
\end{equation}
If one passes to generalized simplices and a vertex becomes ideal, then the corresponding term in
$w_0$ vanishes by continuity.  In particular, for an ideal $6$-simplex one gets the simplified
formula
\begin{equation}
\label{eq:h6-ideal}
\Vol_6(\Delta_{\mathrm{ideal}})
=
\frac{8\pi^3}{15}
\left(
\frac{17}{4}-w_4+\frac12 w_2
\right).
\end{equation}

\subsection*{How the code evaluates the $\Ss^5$ links}

Let $B$ be a $6\times 6$ spherical angle Gram matrix.  The associated spherical vertex Gram matrix is
\begin{equation}
\label{eq:spherical-vertex-gram}
 G_{ab}=\frac{(B^{-1})_{ab}}{\sqrt{(B^{-1})_{aa}(B^{-1})_{bb}}}.
\end{equation}
Choose vertices $v_0,\dots,v_5\in \Ss^5\subset \Rr^6$ with $v_a\cdot v_b=G_{ab}$.
To put the simplex into gnomonic coordinates, solve
\[
G\,\alpha=\mathbf 1,
\qquad
c'=\sum_{a=0}^{5}\alpha_a v_a,
\qquad
c=\frac{c'}{\|c'\|}\in \Ss^5.
\]
Then $c$ lies in the interior hemisphere of the simplex.  If $e_1,\dots,e_5$ is an orthonormal
basis of $c^{\perp}$, the gnomonic coordinates are
\begin{equation}
\label{eq:gnomonic-coords}
 x_a=
 \frac{1}{v_a\cdot c}
 \bigl((v_a\cdot e_1),\dots,(v_a\cdot e_5)\bigr)
 \in \Rr^5.
\end{equation}
Under the inverse gnomonic map
\[
 x\longmapsto \frac{c+\sum_{r=1}^{5}x_r e_r}{\sqrt{1+\|x\|^2}},
\]
the spherical volume element becomes
\begin{equation}
\label{eq:spherical-gnomonic-density}
 d\Vol_{\Ss^5}=
 \frac{dx}{(1+\|x\|^2)^3}.
\end{equation}
Hence the code evaluates each vertex link by
\begin{equation}
\label{eq:s5-qmc}
\Vol_{\Ss^5}(L_i)
=
\int_{\conv(x_0,\dots,x_5)}\frac{dx}{(1+\|x\|^2)^3},
\end{equation}
again using Sobol QMC on the Euclidean simplex in gnomonic coordinates.

\section{Practical remarks and what transfers to polytopes}

\begin{enumerate}[leftmargin=2.2em]
\item For $\Hh^5$, the code's main routine is the Schlaefli path integrator
\eqref{eq:h5-radial}.  A direct $5$-dimensional Klein integral is included as a cross-check.

\item For $\Hh^6$, the code's main routine is the absolute reduction formula
\eqref{eq:h6-main}.  A direct $6$-dimensional Klein integral is again included as a cross-check.
Near the Euclidean or very-small-volume regime, \eqref{eq:h6-main} subtracts several large angle-sum
terms and can suffer noticeable cancellation.  In that regime the direct Klein integral is the more
reliable diagnostic.

\item For general polytopes one may triangulate and sum simplex volumes.  In even dimension the Hopf
face-sum formula also extends directly to a polytope, but from the coding perspective triangulation
is usually the simplest path.
\end{enumerate}

\section{Two numerical smoke tests}

With the defaults in the companion file \texttt{hyperbolic\_5\_6\_simplex\_volume.py}, the following
examples are useful quick checks.

\begin{center}
\begin{tabular}{@{}llll@{}}
\toprule
Space & simplex & main routine & direct Klein cross-check \\
\midrule
$\Hh^5$ & $\conv(0,0.5e_1,\dots,0.5e_5)$ & $3.14018\times 10^{-4}$ & $3.14365\times 10^{-4}$ \\
$\Hh^6$ & $\conv(0,0.7e_1,\dots,0.7e_6)$ & $2.44670\times 10^{-4}$ & $2.44975\times 10^{-4}$ \\
\bottomrule
\end{tabular}
\end{center}

The point of the second row is not that the reduction formula is ``more exact'' numerically, but
rather that its constituent lower-dimensional pieces are now explicit enough to make an implementation
practical.  The direct integral remains the best sanity check.

\begin{thebibliography}{99}

\bibitem{Heckman}
G. J. Heckman,
\emph{The volume of hyperbolic Coxeter polytopes of even dimension},
available at \url{https://www.math.ru.nl/~heckman/Heck_7.pdf}.

\bibitem{KellerhalsZehrt}
R. Kellerhals and T. Zehrt,
\emph{The Gauss-Bonnet Formula for Hyperbolic Manifolds of Finite Volume},
Geom. Dedicata \textbf{84} (2001), 49--62;
preprint available at \url{https://homeweb.unifr.ch/kellerha/pub/GaussBonnet.pdf}.

\bibitem{Kohno}
T. Kohno,
\emph{The volume of a hyperbolic simplex and iterated integrals},
available at \url{https://www.ms.u-tokyo.ac.jp/~kohno/papers/VHI01.pdf}.

\bibitem{Kellerhals5}
R. Kellerhals,
\emph{Volumes in hyperbolic 5-space},
Geom. Funct. Anal. \textbf{5} (1995), 640--667;
preprint available at \url{https://homeweb.unifr.ch/kellerha/pub/gafa.pdf}.

\bibitem{MurakamiSpherical}
J. Murakami,
\emph{Volume formulas for a spherical tetrahedron},
Osaka J. Math. \textbf{49} (2012), 631--638;
preprint available at \url{https://arxiv.org/pdf/1011.2584}.

\bibitem{MY}
J. Murakami and M. Yano,
\emph{On the Volume of a Hyperbolic and Spherical Tetrahedron},
Comm. Anal. Geom. \textbf{13} (2005), 379--400;
preprint available at \url{https://murakami.w.waseda.jp/papers/tetrahedronrev3.pdf}.

\bibitem{Ushijima}
A. Ushijima,
\emph{A volume formula for generalized hyperbolic tetrahedra},
Discrete Comput. Geom. \textbf{31} (2004), 305--314;
preprint available at \url{https://arxiv.org/abs/math/0309216}.

\end{thebibliography}

\end{document}
